\setchapterpreamble[u]{\margintoc}
\chapter{Linear Algebra in Optimisation}
\labch{optimization}

\sources{Strang, \emph{Introduction to Linear Algebra}, 6th ed.\ \cite{strang2023ila},
Chapter~9.}

Optimisation problems that can be solved exactly are, with few exceptions,
quadratic ones. This chapter collects the linear algebra that governs them: the
norms used to measure error, the condition number that predicts difficulty, and
the quadratic form whose minimiser is a linear solve in disguise. The calculus
and the algorithms come in \refch{gradient-methods}.

\section{Norms}
\labsec{norms}

For vectors the three standard choices are
\[
	\lVert\vect{x}\rVert_1=\sum_i|x_i|,
	\qquad
	\lVert\vect{x}\rVert_2=\Bigl(\sum_i x_i^{2}\Bigr)^{1/2},
	\qquad
	\lVert\vect{x}\rVert_\infty=\max_i|x_i| .
\]
All three satisfy the triangle inequality; only \(\lVert\cdot\rVert_2\) comes
from an inner product, which is why projection and least squares are stated in
it.

\marginnote{The choice is not cosmetic. Penalising \(\lVert\cdot\rVert_1\)
drives coefficients exactly to zero and produces sparse solutions; penalising
\(\lVert\cdot\rVert_2\) shrinks them smoothly and keeps them non-zero.}

For matrices, two norms are used constantly:
\[
	\lVert A\rVert_2=\sigma_1,
	\qquad
	\lVert A\rVert_F=\Bigl(\sum_{i,j}a_{ij}^{2}\Bigr)^{1/2}
	=\Bigl(\sum_i\sigma_i^{2}\Bigr)^{1/2},
\]
the largest singular value and the Frobenius norm. Both are supplied by the
decomposition of \refch{svd}; the second identity says the Frobenius norm does
not care which orthonormal basis the matrix is written in.

\section{Conditioning}
\labsec{conditioning}

\begin{definition}
The \emph{condition number} of an invertible \(A\) is
\(\kappa(A)=\sigma_1/\sigma_n\), the ratio of largest to smallest singular value.
\end{definition}

It bounds how much a relative perturbation of \(\vect{b}\) can be amplified in
the solution of \(A\vect{x}=\vect{b}\): a relative change of size \(\varepsilon\)
in the data produces a relative change of at most \(\kappa(A)\varepsilon\) in the
answer. A matrix with \(\kappa=10^{8}\) loses roughly eight digits, whatever
algorithm is used.

\begin{remark}
This is the precise reason \refch{orthogonality} warned against forming
\(A\tp A\). Its singular values are the squares of those of \(A\), so
\(\kappa(A\tp A)=\kappa(A)^{2}\) and the difficulty is squared along with them.
\end{remark}

\section{Quadratic forms}
\labsec{quadratic-forms}

A quadratic function on \(\R^{n}\) can always be written
\[
	f(\vect{x})=\tfrac12\vect{x}\tp A\vect{x}-\vect{b}\tp\vect{x}+c,
\]
with \(A\) symmetric --- any non-symmetric part cancels in the product and can be
discarded without changing \(f\).

\begin{theorem}
If \(A\) is symmetric positive definite then \(f\) has a unique minimiser, and it
is the solution of \(A\vect{x}=\vect{b}\).
\end{theorem}

\begin{proof}
Let \(A\vect{x}^{\star}=\vect{b}\) and write \(\vect{x}=\vect{x}^{\star}+\vect{d}\).
Expanding and using symmetry,
\[
	f(\vect{x})=f(\vect{x}^{\star})+\tfrac12\vect{d}\tp A\vect{d},
\]
and positive definiteness makes the last term strictly positive unless
\(\vect{d}=\vect{0}\).
\end{proof}

\marginnote{Minimisation and solving are therefore the same problem wearing two
hats. Least squares in \refch{orthogonality} was the first instance:
\(\lVert A\vect{x}-\vect{b}\rVert^{2}\) is a quadratic form whose matrix is
\(A\tp A\).}

The level sets of \(f\) are ellipsoids whose axes are the eigenvectors of \(A\),
with lengths proportional to \(1/\sqrt{\lambda_i}\). A large ratio
\(\lambda_{\max}/\lambda_{\min}\) makes them long and thin, which is exactly the
geometry that slows gradient methods down in \refch{gradient-methods}.

\section{Regularisation}
\labsec{regularisation}

When \(A\tp A\) is singular or nearly so, adding a multiple of the identity
repairs it:
\[
	(A\tp A+\delta^{2}I)\hat{\vect{x}}=A\tp\vect{b},
	\qquad \delta>0 .
\]
The matrix on the left is positive definite for every \(\delta>0\), since
\(\vect{x}\tp(A\tp A+\delta^{2}I)\vect{x}
=\lVert A\vect{x}\rVert^{2}+\delta^{2}\lVert\vect{x}\rVert^{2}\). In terms of the
singular value decomposition each \(\sigma_i\) is replaced by
\(\sigma_i/(\sigma_i^{2}+\delta^{2})\), so small singular values are damped
rather than inverted.

\begin{remark}
Statistics calls this ridge regression and deep learning calls it weight decay,
but the linear algebra is identical: trade a little bias for a large reduction
in sensitivity to noise. As \(\delta\to0\) the solution approaches the
pseudoinverse answer of \refsec{pseudoinverse}.
\end{remark}
